% Einheit 2 von 4 – Normalenvektor und Koordinatengleichung (bank/ebenen/e2.jsonl, zusammenbau v0.1)
%% TODO Zeitmarke Einheit 2: Marken-Zeile nicht lesbar
%% TODO Zweigzeile Teil 1: was hier gelernt wird (Satz in Schülersprache aus dem Einheitstitel) – Einheit 2
\einheitenkopf[e2]{Einheit 2 von 4 · Normalenvektor und Koordinatengleichung}
\zweigzeile{Abitur GK}
\setcounter{aufgabe}{11}

%% TODO \verfahren{…}: Verfahrensüberschrift in Sachsprache fehlt in der Bank (2.3 b)
%% TODO Titel: Ich-kann-Satz fehlt in der Bank (Platzhalter: Kettenname) – Nr. 12
%% TODO Anweisung: ein Satz über der ersten Teilaufgabe fehlt in der Bank – Nr. 12
\begin{aufgabe}{Koordinatengleichung bestimmen}
%% TODO \rechenplatz{n}: Schrittzahl des Grundfalls fehlt in der Bank (nur bei fester Schreibform, 2.3 b)
\begin{teile}
\teil Gesucht ist eine Gleichung der Form $ax + by + cz = d$ für die Ebene durch drei gegebene Punkte. Welche Form ist verlangt? \\ \kreuz{Parameterform} \\ \kreuz{Koordinatenform} \\ \kreuz{Normalenform}
\teil Spannvektoren $(1 | 0 | 3)$ und $(0 | 1 | 2)$: ein Normalenvektor $\vec n$ der Ebene? $\vec n = ($ \leerfeld | \leerfeld | \leerfeld $)$
\teil Ein dreieckiges Sonnensegel ist an drei Masten in $A(3 | 2 | 2)$, $B(3 | 5 | 2)$ und $C(0 | 3 | 4)$ befestigt (1 LE = 1 m). Die Ebene des Segels hat eine Gleichung der Form $2x + 3z = d$. Bestimme $d$. \hfill (Abitur 2020 GK) \\ $d = $ \leerfeld
\teil Ein Dachaufbau hat die Form einer Pyramide mit der Grundfläche $A(1 | 1 | 0)$, $B(4 | 1 | 0)$, $C(4 | 4 | 0)$, $D(1 | 7 | 0)$ und der Spitze $S(1 | 1 | 5)$ (1 LE = 1 m). Die Seitenfläche $CDS$ liegt in der Ebene $E$. Bestimme eine Gleichung von $E$ in Koordinatenform (zur Kontrolle: $5x + 5y + 6z = 40$). \hfill (Abitur 2024 GK)
\teil Bestimme eine Koordinatengleichung von $E\colon \vec x = (4 | 1 | 0) + r \cdot (-2 | 4 | 0) + s \cdot (-2 | 0 | 3)$, $r, s \in \mathbb{R}$ (zur Kontrolle: $6x + 3y + 4z = 27$). \hfill (Abitur 2020 GK)
\teil Ein Sonnensegel $W$ ist in $M(3 | 3 | 2)$, $F(0 | 0 | 3)$ und $S(6 | 0 | 0)$ befestigt (1 LE = 1 m). Bestimme eine Koordinatengleichung von $W$ mit dem Ansatz $ax + by + cz = 18$ (zur Kontrolle: $3x - y + 6z = 18$). \hfill (Abitur 2024 GK)
\teil Eine dreieckige Plane über einer Lagerhalle hat die Ecken $U(1 | 1 | 5)$, $V(3 | 9 | 5)$ und $W(9 | 1 | 3)$ (1 LE = 1 m). Gib die Ebene $E$ der Plane in Parameter- und in Koordinatenform an (zur Kontrolle: $4x - y + 16z = 83$). \hfill (Abitur 2020 GK)
\teil Die Geraden $g\colon \vec x = (1 | 3 | 2) + t \cdot (2 | 0 | 3)$ und $h\colon \vec x = (3 | 3 | 5) + k \cdot (1 | 0 | -1)$, $t, k \in \mathbb{R}$, schneiden sich in $(3 | 3 | 5)$. Die Ebene $E$ enthält beide. Gleichung von $E$ in Koordinatenform? \hfill (Abitur 2018 LK)
\teil Für jedes $k > 0$ ist $D_k(0 | 0 | k)$ die Spitze eines Körpers; seine Seitenfläche $BCD_k$ mit $B(9 | 0 | 0)$ und $C(0 | 3 | 0)$ liegt in der Ebene $L_k$. Bestimme eine Koordinatengleichung von $L_k$ (zur Kontrolle: $x + 3y + \frac{9}{k} \cdot z = 9$). \hfill (Abitur 2022 LK)
\end{teile}
\end{aufgabe}

%% TODO \verfahren{…}: Verfahrensüberschrift in Sachsprache fehlt in der Bank (2.3 b)
%% TODO Titel: Ich-kann-Satz fehlt in der Bank (Platzhalter: Kettenname) – Nr. 13
%% TODO Anweisung: ein Satz über der ersten Teilaufgabe fehlt in der Bank – Nr. 13
\begin{aufgabe}{Normalenvektor und Nachweise}
%% TODO \rechenplatz{n}: Schrittzahl des Grundfalls fehlt in der Bank (nur bei fester Schreibform, 2.3 b)
\begin{teile}
\teil Gegeben ist $E\colon \vec x = \vec p + r \cdot \vec u + s \cdot \vec v$ mit Zahlen für die Vektoren. In welcher Form ist $E$ gegeben? \\ \kreuz{Parameterform} \\ \kreuz{Koordinatenform} \\ \kreuz{Normalenform}
\teil Gib zwei Normalenvektoren von $E\colon 4x - 2y + z = 7$ an.
\teil Gegeben ist $F\colon \left(\vec x - (2 | 5 | 1)\right) \cdot (-3 | 1 | 4) = 0$. Gib $F$ in Koordinatenform an. \hfill (Abitur 2021 GK)
\teil Gegeben sind $E(5 | 2 | 3)$, $F(1 | 4 | 3)$ und $S(1 | 1 | 5)$. Bestimme rechnerisch einen Normalenvektor der Ebene $EFS$ und prüfe ihn. \hfill (Abitur 2025 GK)
\teil Eine dreieckige Teilfläche eines Pavillondachs liegt in der Ebene $H$, die die Gerade $g\colon \vec x = (-2 | 2 | 3) + t \cdot (-4 | 2 | -1)$, $t \in \mathbb{R}$, und den Punkt $P(2 | 4 | 2)$ enthält. Weise nach, dass $H$ durch $y + 2z = 8$ beschrieben werden kann. \hfill (Abitur 2017 LK)
\teil Die Ebene $E$ enthält $B(4 | 1 | 0)$, $C(1 | 4 | 0)$ und $S(2 | 2 | 3)$. Begründe ohne $\vec n$ zu berechnen, dass jeder Vektor $\vec n \ne \vec 0$ mit $\vec n \cdot (1 | -1 | 0) = 0$ und $\vec n \cdot (-4 | 2 | 6) = 0$ ein Normalenvektor von $E$ ist. \hfill (Abitur 2025 LK)
\end{teile}
\end{aufgabe}

%% TODO Titel: Ich-kann-Satz fehlt in der Bank (Platzhalter: Kettenname) – Nr. 14
%% TODO Anweisung: ein Satz über der ersten Teilaufgabe fehlt in der Bank – Nr. 14
\begin{aufgabe}{Koordinatengleichung bestimmen – Fehler finden, Begründen, Anwendung, Darstellungswechsel}
\begin{teile}
\teil Nele bestimmt die Ebene durch $A(2 | 1 | 3)$, $B(4 | 0 | 3)$ und $C(2 | 4 | 1)$. Sie findet richtig $\vec n = (1 | 2 | 3)$ und schreibt: „Ich setze den Ursprung ein: $E\colon x + 2y + 3z = 0$.“ Finde den Fehler.
\teil Begründe, warum mit einem Normalenvektor $\vec n$ auch jedes Vielfache $k \cdot \vec n$ ($k \ne 0$) ein Normalenvektor derselben Ebene ist.
\teil Eine Dachfläche geht durch $A(0 | 0 | 3)$, $B(8 | 0 | 3)$ und $C(8 | 5 | 6)$ (1 LE = 1 m, die $xy$-Ebene ist der Boden). Wie hoch liegt das Dach über dem Bodenpunkt $(4 | 2 | 0)$?
\teil Schreibe $E\colon 2x - 3y + z = 5$ in Normalenform.
\end{teile}
\end{aufgabe}
